%!TEX program = xelatex
\documentclass[12pt, b5paper]{article}
\usepackage{ctex}

\usepackage{geometry}
\geometry{b5paper,left=2cm,right=1cm,top=2.5cm,bottom=2.5cm}
\usepackage{chemformula} %化学公式宏包
\usepackage[version=4]{mhchem} %化学公式宏包
\usepackage{siunitx} %数字，单位宏包
\sisetup{
	separate-uncertainty = true,
	inter-unit-product = \ensuremath{{}\cdot{}}
} %单位间加小圆点
\usepackage{tikz} %Tikz绘图包
\usetikzlibrary{cd}
\usetikzlibrary{chains}
\usetikzlibrary{arrows}
\usetikzlibrary{arrows.meta} %画箭头
\usetikzlibrary{commutative-diagrams} %交换图
\usetikzlibrary{positioning,automata}
\usepackage[all]{xy}
\usepackage{booktabs} %三线表
\usepackage{multirow} %合并单元格
\usepackage{makecell} %表格内强制换行
\usepackage{array}
\usepackage{booktabs} %控制表格行高
\usepackage{colortbl}  %彩色表格需要加载的宏包
\usepackage{amsmath}
\usepackage{unicode-math}
\usepackage{fontspec} %字体
\newcommand*{\cred}{\textcolor{red}}

\setmainfont{Times New Roman}
\setmathfont{XITS Math}

\setlength{\parindent}{0pt} %无缩进
\usepackage{setspace}
\usepackage{fancyhdr} %设置页码
\usepackage{lastpage} %获取总页数
\pagestyle{fancy} %fancyhdr宏包新增的页面风格
\renewcommand\headrulewidth{0pt} %取消页眉中的装饰分割线
\fancyhf{}
\cfoot{\footnotesize 第 \thepage\ 页(共 \pageref{LastPage}页)}%当前页 of 总页数

\begin{document}
\begin{spacing}{1.625}

\centerline{{\huge 河北农业大学课程考试试卷}}

\centerline{\large 2022-2023学年第\underline{\ 1\ }学期}

考试科目：\underline{物理化学与胶体化学} \quad 姓名：\underline{ \qquad \qquad} \quad 学号：\underline{ \qquad \qquad \qquad}

学院：\underline{理工学院} \quad 专业班级：\underline{\qquad \qquad \qquad \qquad} \quad 卷别：\cred{B} 考核方式：\cred{开卷}

\bigskip


1. 简答题(15分)

生活中的热力学现象无处不在，热力学现象的本质和原理亦来自生活。请列举至少3种生活中的热力学现象，并说明其中包含了哪些热力学原理？

\bigskip

\bigskip

2. 简答题(15分)

化学热力学是物理化学的一个分支，它主要研究物质系统在各种条件下的物理和化学变化中所伴随着的能量变化，从而对化学反应的方向和进行的程度作出准确的判断。如何判断化学反应的方向和进行的程度？

\bigskip

\bigskip

3. 简答题(10分)

根据相律解释以下事实：
\ch{CaCO3}(s)在高温下分解为CaO(s)和\ch{CO2}(g)：

(1)若在定压\ch{CO2}气体中将CaCO3(s)加热，实验证明加热过程中，在一定温度范围内\ch{CaCO3}不会分解；

(2)若保持\ch{CO2}压力恒定，实验证明只有一个温度能使\ch{CaCO3}和CaO混合物不发生变化。


\bigskip

\bigskip

4. 主观题(20分)

计算1mol 298K、101.325kPa的过冷水蒸气变成同温同压下的液态水的$\Delta G$ 。该温度下液态水的饱和蒸汽压为3.168kPa，摩尔体积$V_\mathrm{m} = 0.018 \si{dm^3.mol^{-1}}$。

\bigskip

\bigskip


5. 主观题(20分)

试计算在珠穆朗玛峰顶上的沸水中煮熟一个鸡蛋需要多长时间。

己知组成蛋的蛋白质的热变作用为一级反应，其活化能约为\SI{85}{kJ.mol^{-1}}，在$p^\ominus$下的沸水中煮熟一个鸡蛋需要\SI{10}{min}。水的摩尔蒸发热$\Delta H_\mathrm{m} = \SI{41.05}{kJ.mol^{-1}}$。珠穆朗玛峰顶大气压力$p=\SI{31}{kPa}$。

\bigskip

\bigskip

6. 主观题(20分)

某反应\ch{A + B -> P}，设计实验，测定其反应速率方程。


\end{spacing}
\end{document}
